%%%%%%%%%%%%%%%%%%%%%%%%%
%%%%%%%%%%%%%%%%%%%%%%%%%
% !TeX encoding = UTF-8
% !Mode:: "TeX:UTF-8"
\documentclass[twoside]{thesis-cls-Miya} %

\begin{document}
	
\title{基于预训练模型的用户评论多标签}
\titleB{智能审核系统的设计与实现}
\englishtitleA{A Multi-Label Intelligent Audit System}
\englishtitleB{for User Comments on Pre-Trained Models}

\author{江昊烨}

\kongF{2214100806}
\kongG{22网络空间安全2班}
\kongH{刘玲君 讲师}
\kongI{2026年3月18日}

\maketitle

\frontmatter

\begin{abstract}
近年来，随着互联网的不断发展，各大互联网平台中的用户生成内容（UGC）数量急剧增加。然而，评论区中也随之出现了大量色情、辱骂及各类歧视性言论。这些不良信息既影响了普通用户的浏览体验，也给平台带来了巨大的合规压力。传统的纯人工审核成本过高且速度较慢，很难应对海量数据的实时审核；而普通的敏感关键词过滤又很容易被网民的拼音、谐音缩写等变体绕过，导致漏检率偏高。为了解决上述痛点，本课题基于 RoBERTa 预训练语言模型，设计并开发了一套针对用户评论的多标签智能审核系统。

在数据准备阶段，本文参考了开源项目 Chinese-offensive-language-detect 的数据构造思路。考虑到中文网络评论往往带有较强的口语化和变体特征，本课题对开源项目 Chinese-offensive-language-detect的原始数据集进行了多轮清洗、去重和标准化，最终整理出一个包含 53509 条样本的中文多标签数据集。该数据集涵盖了色情、辱骂、地域歧视、性别歧视、种族歧视以及职业歧视共 6 类风险标签。在算法设计上，本课题选择了针对中文语境优化的 RoBERTa-wwm-ext 作为基础特征提取器。因为一条评论可能会同时包含多种违规行为，本课题采用 Sigmoid 函数独立输出各项违规概率，并配合 BCEWithLogitsLoss 损失函数完成了模型的微调。实验测试显示，微调后的模型 Micro-F1 达到了 0.9921，Macro-F1 达到了 0.9919，分类精度较为理想。

除此之外，由于真实的业务审核通常遵循“宁可误报、不可漏报”的原则，本课题并没有完全依赖单一的黑盒模型，系统在模型推理的基础上，加入了一套传统的规则引擎作为兜底。这种规则与模型相融合的机制，能有效减少高危违规内容的漏检情况。在此基础上，本课题利用 Streamlit 框架搭建了一个可视化的前端交互系统。该原型不仅支持单条评论检测和批量表格数据处理，还可以动态调节各类标签的判定阈值，并能详细记录每次预测的日志与触发的规则。整体而言，本系统实现了一套“机器初筛+人工复核”的工作闭环，在有效分担人工审核压力的同时，也让审核过程更加透明和可追溯，具有较好的工程落地价值。

\keywords{预训练模型；多标签分类；内容审核；RoBERTa；Streamlit；人机协同；网络空间安全}
\end{abstract}

\begin{englishabstract}
In recent years, with the continuous development of the internet, the volume of user-generated content (UGC) on major online platforms has increased dramatically. Consequently, a significant amount of inappropriate content, including pornography, insults, and various forms of discriminatory speech, has emerged in comment sections. This undesirable information not only diminishes the browsing experience for general users but also imposes substantial compliance pressure on platforms. Traditional purely manual moderation is too costly and slow for real-time review of massive data streams, while simple sensitive keyword filtering is easily bypassed by netizens using variations like pinyin, homophones, and abbreviations, resulting in a high rate of missed detections.

To address these challenges, this study designs and develops a multi-label intelligent comment moderation system based on the RoBERTa pre-trained language model. During the data preparation phase, this research drew upon the data construction methodology of the open-source project "Chinese-offensive-language-detect." Recognizing that Chinese online comments often exhibit strong colloquial characteristics and variant forms, we performed multiple rounds of cleaning, deduplication, and standardization on the original dataset from the Chinese-offensive-language-detect project. This process resulted in a Chinese multi-label dataset comprising 53,509 samples. The dataset covers six categories of risk labels: pornography, insults, regional discrimination, gender discrimination, racial discrimination, and occupational discrimination. For the algorithmic design, we selected RoBERTa-wwm-ext, which is optimized for the Chinese context, as the core feature extractor. Given that a single comment may contain multiple types of violations, we employed the Sigmoid function to independently output the probability for each violation type and fine-tuned the model using the BCEWithLogitsLoss function. Experimental results demonstrate that the fine-tuned model achieves a Micro-F1 score of 0.9921 and a Macro-F1 score of 0.9919, indicating highly satisfactory classification accuracy.

Furthermore, recognizing that real-world business moderation typically adheres to the principle of "better false positives than missed detections," this study does not rely solely on a single black-box model. The system integrates a traditional rule engine as a fallback mechanism on top of the model's inference. This hybrid approach, combining rules with the model, effectively reduces the rate of missed detections for high-risk prohibited content. Building on this, we utilized the Streamlit framework to develop a visual front-end interactive system. This prototype not only supports single comment detection and batch processing of tabular data but also allows for dynamic adjustment of decision thresholds for each label and maintains detailed logs of each prediction and the rules triggered. Overall, this system establishes a closed-loop workflow of "machine pre-screening + human review." While effectively alleviating the pressure of manual moderation, it also renders the moderation process more transparent and auditable, demonstrating significant potential for practical engineering application.

\englishkeywords{Pre-trained Models; Multi-label Classification; Content Moderation; RoBERTa; Streamlit; Human-Machine Collaboration; Cybersecurity}
\end{englishabstract}

\tableofcontents

\mainmatter

\chapter{绪论} 

\section{研究背景及意义}
如今，各类社交平台和社区论坛已经成为网民交流的主要场所。诸如评论、弹幕等形式的用户生成内容（UGC）每天都在大量产生，并且传播速度极快。高质量的评论能够活跃社区氛围；但与此同时，由于发布门槛较低，评论区也经常会混杂着色情、辱骂以及各种歧视性的不良言论。这些违规内容不仅会影响普通用户的浏览体验，也给平台的日常监管和合规工作带来了巨大的压力。

在管理这些内容时，目前大部分平台主要依靠“人工审核+关键词过滤+普通机器学习”的组合方式。但是，这种传统做法的缺陷越来越明显：人工审核虽然准确但成本太高，很难应付每秒钟成千上万条的实时数据；普通的敏感词过滤又很容易被拼音、谐音缩写等变体绕过，导致漏检率一直居高不下。而早期的普通机器学习模型（如TextCNN），因为对上下文语义的理解不够深，在面对复杂的网络用语或需要判断多重违规标签时，效果往往不太理想。

基于上述问题，开发一套智能化的多标签内容审核系统显得非常有必要。从实际应用的角度来看，利用深度学习模型对海量评论进行自动化初筛，可以帮审核人员分担掉绝大部分工作量。同时，相比于传统的敏感词匹配，预训练模型能够根据上下文理解一些比较隐晦的骂人方式，从而减少漏检的风险。除此之外，如果我们能将模型预测结果记录下来，并搭配明确的规则提示，还能让整个审核过程更加透明和可追溯，为网络平台的合规管理提供技术支持。

\section{国内外发展现状分析}

\subsection{国内发展现状}
国内的互联网内容审核技术发展非常快，主要是为了应对国内庞大的网民基数和复杂的中文网络环境。整体来看，国内的技术路线经历了一个从纯人工向大模型辅助转变的过程。

在早期阶段，平台主要雇佣大量审核专员，配合最基础的敏感词库来进行文本过滤。但随着短视频和社交媒体的爆发，这种低效的方法很快就跟不上了。后来，行业内开始引入传统的机器学习和早期的深度学习算法（如 TextCNN、LSTM 等），这让审核系统具备了一定的语义分析能力，能够自动拦下一些简单的违规内容。到了现阶段，像 BERT、RoBERTa-wwm-ext 等针对中文优化过的预训练模型已经非常普及。它们能够很好地结合上下文来判断某句话是不是“阴阳怪气”或者在骂人，大大提高了筛查的准确率。

目前，国内审核行业主要有几个明显的趋势：首先是标签划分越来越细。以前只是判断“违规”或“不违规”，现在往往需要明确指出到底是涉黄、辱骂还是地域歧视。其次是多模态技术的落地，很多平台已经把文本、图片甚至视频的审核整合到了一起。最后，也是最贴近真实业务的一点，就是人机协同审核已经成为标配，机器负责处理 90\% 的明显违规或安全内容，剩下那些模棱两可的再交给人去复核。不过，目前对于一些小众的网络黑话或者方言，模型的识别能力依然有待提高。

\subsection{国外发展现状}
国外在不良内容检测方面的研究起步相对较早。受文化背景影响，他们的研究重点更多放在仇恨言论（Hate Speech）、种族歧视和暴力内容的识别上。

在技术演进上，国外也经历了类似的过程。早些年大家都在用 TF-IDF 提取特征，然后扔给 SVM 或逻辑回归跑分类。之后，随着 Transformer 架构的提出，RoBERTa、DeBERTa 等预训练模型迅速成为了海外大型互联网企业（如 Facebook、YouTube）的审核标配。近两年，由于大语言模型（如 Llama、ChatGPT）的爆发，国外的研究热点开始转向如何利用大模型的推理能力，在没有太多标注数据的情况下（少样本或零样本），去识别新出现的违规词汇，甚至让模型自己解释为什么违规。

在学术界，国外发布了许多经典的数据集，比如 Toxic Comment Classification 和 HateEval，这极大方便了研究者进行模型对比。相比于国内，国外的研究者可能会花更多的心思在“模型公平性”上，也就是想办法避免模型对某些特定人群产生误判。虽然国内外关注的违规重点不太一样，但在技术路线上殊途同归，基于预训练模型的多标签分类目前都是主流解法。这也是本文选择该技术方向的重要原因。

\section{研究内容及论文组织结构}

\subsection{主要研究内容}

为了解决中文用户评论中存在的违规问题，本课题结合深度学习与人机协同的思路，设计并实现了一套多标签审核系统。具体的研究内容包括以下几个方面：

1. 构建多标签审核数据集：参考开源项目 Chinese-offensive-language-detect 的数据结构，结合中文网络环境的特点，本课题清洗并整理了一份包含 53509 条样本的数据集。该数据集主要涵盖了色情、辱骂以及地域、性别、种族、职业等四类歧视标签。

2. 多标签分类模型的训练与选型：在代码实现层面，我们对比了 MacBERT 和 RoBERTa-wwm-ext 两种预训练模型的表现。针对多标签问题，本课题使用了 Sigmoid 激活函数和 BCEWithLogitsLoss 损失函数对模型进行微调，以此来提升模型对隐晦违规文本的分类精度。

3. 设计规则融合的兜底策略：考虑到真实的业务审核通常要求“宁可误报，不可漏报”，本课题没有单纯依靠模型去盲猜。相反，本课题在模型输出概率的基础上，加入了一套传统的规则匹配引擎（如敏感词典），将两者得分融合，以此来进一步降低漏检率。

4. 开发可视化审核系统：最后，本课题使用 Streamlit 框架开发了一个轻量级的前端交互系统。它不仅能让审核员直接测试单条文本，还支持上传表格进行批量检测，并允许管理员在后台动态调整各个违规标签的判定阈值。

\subsection{论文组织结构}

为了清晰地展示上述研究过程，本文共分为六个章节，具体安排如下：

第一章是绪论。主要介绍了本课题的背景和开发这套审核系统的实际意义，并梳理了国内外在内容审核领域的技术发展现状。

第二章介绍了相关的理论基础。重点讲解了自然语言处理中的多标签分类概念，以及本课题用到的 Transformer 架构和预训练语言模型（特别是 RoBERTa-wwm-ext）的底层原理。

第三章详细说明了模型的结构和我们在工程上的改进策略。这一章解释了数据是如何输入到预训练编码器中的，分类头是如何设计的，以及我们是如何通过公式把规则得分与模型概率融合在一起的。

第四章展示了模型的训练过程与结果分析。包括数据集的清洗方式、实验环境和超参数的设定，最后通过图表展示了模型对比的结果，并对出现误判和漏判的样本进行了误差分析。

第五章介绍了智能审核系统的工程实现。包括整个软件的架构设计、各个功能模块（如单条预测、批量处理、动态调节阈值）的实现代码逻辑，以及系统的交互流程。

第六章是总结与展望。对我本人在本次毕业设计中完成的工作进行了回顾，指出了现有系统还可以继续优化的短板，并对未来的技术迭代方向进行了设想。


\chapter{相关理论基础}

\section{自然语言处理与多标签文本分类}

自然语言处理（NLP）是人工智能的一个重要分支，主要目的是让计算机能够理解和处理人类的语言。在本课题中，我们要做的用户评论违规审核，底层就是依赖 NLP 技术来实现的。文本分类是 NLP 里面非常基础的一项工作，简单来说，就是把一段没有固定格式的自然语言，自动归类到我们提前设定好的标签里。这项技术在情感分析、垃圾邮件拦截以及我们关注的内容审核中都用得非常多。

根据标签的数量和相互关系，文本分类一般可以分成单标签分类和多标签分类两种。本课题要解决的评论审核问题，就是一个典型的\textbf{多标签文本分类任务}。以前的单标签分类，各个类别之间往往是互斥的（比如一句话要么是正面，要么是负面）。但在真实的评论区里，情况要复杂得多。一句违规的话，可能既包含了“辱骂”，又带有“性别歧视”。这就要求本课题的模型不能只输出一个非黑即白的结果，而是要在理解整段话的意思后，同时判断出它是否触发了多个不同的风险标签。

\section{Transformer 架构}
提到现在的自然语言处理技术，就不得不提谷歌在 2017 年提出的 Transformer 架构。在这之前，大家处理文本序列主要用的是 RNN 或者 LSTM 这些循环神经网络，但这类网络在处理长句子时容易遗忘前面的信息，而且很难进行并行计算，导致训练速度比较慢。Transformer 打破了这个常规，它直接放弃了传统的循环结构，完全靠“自注意力机制（Self-Attention）”来理解句子中各个词语前后的关系。

自注意力机制的作用在于，当模型在处理句子里的某一个词时，它可以同时把注意力分配给这句话里的其他所有词。这样一来，模型就能很好地结合上下文来理解那些比较隐晦的网络用语。原本的 Transformer 是由编码器（Encoder）和解码器（Decoder）两部分组成的。不过对于本文的文本分类任务来说，我们不需要让模型去“生成”文字，主要是利用它负责提取特征的编码器部分。本课题选用的 RoBERTa-wwm-ext 模型，也就是在 Transformer 编码器的基础上发展而来的。

\section{预训练语言模型基础}
预训练语言模型的核心思想可以概括为“先通用学习，再特定微调”。也就是先让模型在海量的普通网页文本上进行无监督学习，掌握人类语言的基本规律和词汇关系；然后再用我们自己手里少量的标注数据（比如本课题的违规评论数据集）对它进行微调。这种方法不用再像以前那样人工去设计各种复杂的特征规则，而且准确率也高得多，现在已经成为了 NLP 开发的常规操作。

在这其中，BERT 是比较具有代表性的一个双向预训练模型。它在预训练阶段主要做了两件事：一个是“掩码语言建模（MLM）”，也就是把句子里的词遮住一部分让模型去猜；另一个是“下一句预测（NSP）”。这让它在很多下游任务上都表现得很好。但是原版的 BERT 在处理中文时，一般是按单个汉字来切分的，这其实不太符合中文以“词”为基本表意单位的习惯。

为了解决这个问题，我们在项目中选用了针对本文优化过的 RoBERTa-wwm-ext 模型。相比于原版 BERT，它去掉了不太实用的 NSP 任务，并且引入了“全词掩码（Whole Word Masking, WWM）”技术。具体来说，当它在做填空题时，遮住的不再是单独的汉字，而是一个完整的中文词语。这就迫使模型必须去理解整个词汇的含义，而不是光靠字面去猜。这种机制使得它在理解中文网络评论时更加准确，因此我们把它作为本次多标签审核系统的核心基座。


\chapter{模型结构及改进策略}

\section{总体网络架构与处理流程}
为了能够同时识别出评论中的六类不同违规标签，本课题在设计分类模型时，搭建了一个“输入层 $\rightarrow$ 预训练特征提取 $\rightarrow$ 多标签输出”的三层结构。这套架构相对比较直接，主要是为了在保证分类准确率的同时，也能满足服务器实时推理的速度要求。

在具体的代码实现中，一条评论经过模型的完整流程如下：

\begin{enumerate}
	\def\labelenumi{\arabic{enumi}.}
	\item
	\textbf{文本预处理与分词（Tokenization）}：当一段用户评论进入系统后，程序会调用和 RoBERTa-wwm-ext 配套的 WordPiece 算法对句子进行切分。同时，程序会自动在句子的开头加上 \texttt{{[}CLS{]}} 标识符（用来汇总全局信息），在句尾加上 \texttt{{[}SEP{]}} 标识符（用来标记句子结束）。
	\item
	\textbf{特征嵌入（Embedding）}：切分好的字符序列不能直接输入模型，需要先转换成稠密向量。我们把每个词的词向量、它在句子中的位置向量以及分段向量相加，以此来保留文字本身的含义和语序信息。
	\item
	\textbf{深度特征提取}：接着，这些向量会被送入一个由 12 层 Transformer Encoder 堆叠而成的网络中。模型通过多头自注意力机制（Multi-Head Attention），去计算这句话里面不同词语之间的相互联系。
	\item
	\textbf{多标签分类映射}：经过层层计算后，我们把带有全句综合语义的 \texttt{{[}CLS{]}} 向量单独提取出来，把它送入一个全连接层。配合 Sigmoid 激活函数，模型最终会吐出 6 个在 0 到 1 之间的小数，分别代表触发各项违规标签的概率。
\end{enumerate}

\section{预训练编码器特征提取机制}
在特征提取这部分，本课题选用的 RoBERTa-wwm-ext 主要依靠多头自注意力机制（Multi-Head Self-Attention）来工作。具体来说，当模型在审视句子里的某一个词时，会通过三个不同的权重矩阵 \(W^Q, W^K, W^V\)，算出三个向量：查询向量（Query）、键向量（Key）和值向量（Value）。

系统在计算注意力权重时的公式如下：
\begin{equation}\text{Attention}(Q, K, V) = \text{softmax}\left(\frac{QK^T}{\sqrt{d_k}}\right)V \
\end{equation}
公式里的 \(d_k\) 代表键向量的维度。我们在点积之后除以 \(\sqrt{d_k}\)，主要是为了防止数字变得太大，导致后面算 Softmax 时梯度消失、模型无法更新参数。因为 RoBERTa-wwm-ext 配置了 12 个注意力头，这就相当于让模型从 12 个不同的角度（比如有的关注语法、有的关注情感词、有的关注指代关系）去拆解这句评论。这也是为什么该模型在对付那些“阴阳怪气”或者拐弯抹角的歧视性言论时，比传统算法要准得多的原因。

\section{多标签分类头设计}
当数据跑完整个编码器后，我们拿到了 \([\text{CLS}]\) 对应的隐藏层向量 \(h \in \mathbb{R}^d\)（在本文中这个维度 d 为 768）。考虑到我们手头的训练数据相对有限，如果直接往下传，模型很容易“死记硬背”导致过拟合。所以在接入分类器之前，我们先加入了一个 Dropout 层，让一部分神经元随机“休眠”。随后，我们用一个全连接层（Linear Layer），把原本 768 维的特征压缩到 6 维，正好对应本课题需要检测的 6 个违规标签。

在选择输出层的激活函数时，我们特意避开了多分类任务里最常用的 Softmax。因为 Softmax 会把所有类别的概率加起来凑成 1，它本质上是一个“单选题”，不同标签之间会互相排斥。但在实际的评论审核场景里，这是个“多项判断题”，一句话完全可以既是色情又是辱骂。因此，我们选用了\textbf{Sigmoid激活函数}。它可以独立地把每个维度的输出都压缩到 0 到 1 之间，计算公式如下：
\begin{equation}p_k = \frac{1}{1 + e^{-z_k}}
\end{equation}
公式里 \(z_k\) 就是全连接层算出来的第 k 个标签的原始分值，而 \(p_k\) 就是最终转化后的违规概率。与之对应的，我们在训练时使用了 BCEWithLogitsLoss（二元交叉熵损失）。这个函数直接把 Sigmoid 计算和交叉熵损失合二为一，不仅计算起来更稳定，也能在一定程度上缓解某些少数违规标签（如种族歧视）数据量太少带来的影响。

\section{面向召回优先的规则融合策略}
深度学习虽然在理解复杂句意上表现很好，但它毕竟是个“黑盒”。在实际测试中我们发现，因为语料分布的问题，模型有时候反而会漏掉一些非常直白的违规内容（比如直接发乱码网址、留各种变体微信号、或者直接甩极端辱骂词）。而真实的审核业务中，通常要求“宁可错杀一千，绝不能放过一个违规的”。

为了堵住深度学习模型的这个漏洞，我们在代码里还实现了一套\textbf{“模型概率 + 启发式规则”}的融合引擎。这套规则主要做三件事：
\begin{enumerate}
	\def\labelenumi{\arabic{enumi}.}
	\item
	\textbf{正则匹配}：专门抓取各种变形的联系方式（比如“加v”、“企鹅号”）以及违规 URL。
	\item
	\textbf{敏感词库匹配}：利用 Aho-Corasick 多模式匹配算法，高速扫描文本里有没有包含被列入“黑名单”的绝对违规词。
	\item
	\textbf{刷屏检测}：判断文本是不是在短时间内疯狂重复几个毫无意义的字符。
\end{enumerate}

在实际打分时，假设模型给第 \(k\) 个标签算出的概率是 \(p_k\)，如果我们自己写的规则引擎抓到了问题，就会给出一个加成惩罚项 \(r_k\)（比如只要匹配到高危脏话，\(r_{\text{abuse}}\) 就直接加上 0.5；没抓到就是 0）。系统最终的风险概率计算如下：
\begin{equation}p_k^{\text{fused}}=\min(1,\; p_k + r_k)
\end{equation}
在部署时，我们可以设置一个全局的报警阈值 \(\tau\)。只要有一个标签融合后的概率 \(p_k^{\text{fused}} \ge \tau\)，系统就会自动把这条评论标红，并推送到后台让人工去复看。通过加上这层传统规则的兜底，我们成功弥补了纯预训练模型在某些低级错误上的不可控性。

\chapter{模型训练与结果分析}

\section{数据集构建与预处理}
\subsection{数据集来源与构造思路}
为了训练出能够应对复杂网络环境的审核模型，我们在选择数据时参考了开源项目 Chinese-offensive-language-detect。这个项目的数据主要是通过大语言模型辅助生成的：作者先使用 LLm 生成每个主题的关键词，然后再围绕这些关键词以及模型越狱技术，成对地生成含有害信息的文本和处于正常语境下的安全文本。这种构造方法产生了很多极具迷惑性的“对抗样本”，能够很好地防止模型在训练时“投机取巧”只记表面特征，从而更贴近真实的审核需求。

原始数据一共分为三个子集：色情违规集（SexHarmSet）、辱骂违规集（AbuseSet）和歧视类违规集（BiasSet）。这三大子集刚好覆盖了本课题需要检测的六大类违规标签，由于其贴近真实的社交平台评论风格，我们直接将其作为本次实验的基础语料。

\subsection{数据清洗与整合}
因为这批数据是由多个不同的子集拼凑而成的，在正式开始训练前，我们必须对其进行清洗。首先，我们把各个 CSV 文件的表头统一，把具体的评论文字和对应的标签字段拼接在一张大表里。接着，我们写脚本剔除了表里的空行、乱码以及完全重复的冗余数据。在标签处理方面，我们把六个违规类别全部转换成了标准的二分类形式，如果文本触发了某个标签，对应字段的值就记为 1，否则记为 0。

\subsection{数据集统计与样本划分}
经过上述清洗步骤后，我们最终得到了一份包含 53509 条有效记录的数据集。这其中存在一个很明显的“偏科”现象：色情和辱骂这两类比较常见的违规内容，正样本比例都接近 50\%；但是像地域歧视、性别歧视这类标签，正样本比例只有不到 6\%。虽然这会导致数据分布不平衡，但这恰恰反映了真实网络环境中的数据分布规律。

在对这 5万多条评论进行统计时，我们发现 95\% 以上的句子长度都在 10 到 120 个汉字之间，只有极个别恶意刷屏的文本超过了 200 个字。根据这个统计结果，我们在调用分词器时，把最大截断长度（\texttt{max\_length}）设定为了 150。也就是说，不足 150 字的就用 \texttt{{[}PAD{]}} 补齐，超过 150 字的就直接切断。这样做既保留了绝大部分有用的信息，又避免了显存被浪费在无意义的空白符上。

最后，为了保证模型评估的客观性，我们设定了一个固定的随机种子，把数据集按照 8:1:1 的比例切分成了训练集、验证集和测试集。这三份数据严格隔离，坚决避免数据泄露。

\section{实验环境与超参数配置}
本课题所有模型训练都是在 Linux 环境下完成的，主要使用的硬件是一张显存为32G的 NVIDIA Tesla V100 GPU。开发语言为 Python 3.13，底层的深度学习库用的是 PyTorch 2.6.0，并配合 HuggingFace 的 Transformers 4.57.5 库来加载预训练模型。

在调参阶段，考虑到硬件显存大小和数据特征，我们最终确定的训练配置如下：
\begin{itemize}
	\item
	序列长度：如前文所述设为 150。
	\item
	批大小（batch\_size）：设为 16，这个数值在保证模型梯度稳定更新的同时，不会让 V100 显存溢出。
	\item
	初始学习率（learning\_rate）：设为 2e-5。因为预训练模型已经具备了很好的通用知识，学习率不能调得太大，否则会把模型原有的参数权重破坏掉。
	\item
	训练轮数（epoch）：设为 3，经验表明对于 BERT 类的微调，跑 3 到 4 轮就足以收敛，再多容易发生过拟合。
	\item
	损失函数与优化器：选用 BCEWithLogitsLoss 来计算多标签交叉熵，优化器使用 AdamW。
\end{itemize}

除了以上基础配置，我们还在代码里加入了一套\textbf{“学习率预热（Warmup）+ 权重衰减”}的机制。在训练刚开始的前 10\% 步数里，学习率会从 0 慢慢爬升到 2e-5，让模型先平稳地适应一下我们的任务数据；到了中后期，学习率再慢慢降下来。同时，把 AdamW 的权重衰减系数设为 0.01，用来限制模型参数长得太大。

\section{模型训练过程分析}
我们在训练的 3 个 Epoch 里，设置了每跑 500 个 Step 就记录一次 Loss 变化，并且去验证集上算一下 Macro-F1 的分数。

从训练完成后生成的后台日志文件和曲线来看，模型在第 1 个 Epoch 快结束的时候学得最快，此时 Loss 下降得很陡，验证集的 Macro-F1 分数也迅速超过了 0.95。等跑到第 2 个 Epoch 的中后期，Loss 就不怎么动了，甚至还会稍微反弹。为了拿到真正“最好”的那个模型，我们写了一个\textbf{早停（Early Stopping）}的逻辑：系统会一直盯着验证集的 Macro-F1 分数，如果连续好几个监测点这个分数都不涨了，程序就会自动把当前表现最好的权重保存下来，然后提前结束训练。

\section{模型评估指标}
如前文所述，我们的数据集里存在严重的“偏科”（长尾分布），所以我们没有用普通的“准确率（Accuracy）”来评估模型。否则模型哪怕闭着眼睛把所有歧视类标签都猜成“正常”，它的整体准确率看起来依然很高。

所以，我们选用了 Micro-F1 和 Macro-F1 两个指标。Micro-F1 是把所有类别的对和错混在一起算总账，主要看模型整体的平均实力。而 Macro-F1 则是先给每个具体的违规标签算一个独立的分数，最后再取平均。这种算法“人人平等”，模型只要在那些数量很少的标签（比如职业歧视）上表现不好，整体的 Macro-F1 就会被严重拉低。因此，我们把验证集上的 Macro-F1 当成了最终挑选模型的金标准。

\section{实验结果与误差分析}
\subsection{模型对比实验结果}
在同样的软硬件环境和超参数下，我们把目前中文圈最火的两个开源预训练模型——MacBERT 和 RoBERTa-wwm-ext 放在一起跑了个比赛。具体的分数对比如下表：

表~\ref{tableLabel}不同预训练模型性能对比

\begin{table}[!htb]
	\centering
	\caption{不同预训练模型性能对比}
	\label{tableLabel}
	\renewcommand{\arraystretch}{1.2}
	\setlength{\tabcolsep}{1ex}
	\begin{tabular}{l c p{0.3\textwidth}}
		\toprule
		模型名称 & Micro-F1 & Macro-F1\\
		\midrule
		MacBERT & 0.9876& 0.9862 \\
		RoBERTa-wwm-ext & 0.9921 & 0.9919 \\
		\bottomrule 
	\end{tabular}
\end{table}

结果很明显，RoBERTa-wwm-ext 的两项指标都赢了。特别是它在 Macro-F1 上领先较多，说明它非常擅长抓那些数量很少的歧视类样本。我们分析，这主要归功于它独有的“全词掩码（WWM）”技术——因为很多地域黑、职业黑的话，都非常依赖特定的名词组合，而 RoBERTa 在预训练时就已经强迫自己学会了整词边界，所以它在这方面比 MacBERT 更聪明。这也是我们最终敲定它作为整个审核系统底座的原因。

\subsection{规则融合策略消融实验}
为了证明我们在第三章里自己加的那套“兜底规则”不是多此一举，我们在测试集上专门做了一次对比（消融实验）。

\begin{itemize}
	\item
	\textbf{如果不加规则（纯模型裸奔）}：当测试集里出现一些拼音缩写（比如“nt”、“sb”）或者故意错别字的时候，模型因为在预训练阶段没怎么见过这些“黑话”，算出来的违规概率往往低于 0.5，这就造成了漏报。
	\item
	\textbf{模型+规则一起上}：当我们开启了规则引擎后，只要这些脏词碰到了黑名单词典，系统就会强制给对应的概率加上一个惩罚项 $r_k$。这样一来，它最终的分数就突破了报警红线，被我们成功拦截了下来。
\end{itemize}

虽然加上这套规则后，有时候会不小心误伤好人（比如某人发的是正常的网址），但整体的漏报率（FNR）下降了约 1.2\%。在真实的评论审核环境里，宁可让审核员多看几条安全的评论，也绝对不能让一条高危违规信息流出去。从这个业务逻辑来看，我们加的这套双重保险是非常有价值的。

\subsection{误差成因分析}
在跑完所有测试后，我们把那些模型判错的样本单独拎出来复盘。

我们发现，漏报（FN）主要出在那些太冷门的方言骂人话上，因为训练集里几乎没有这方面的数据，模型自然学不会。而误报（FP）通常是因为有些正常新闻或者科普文章里带了敏感词（比如“裸体雕塑”）；或者是这句话本身就在多种违规的边缘疯狂试探，导致模型自己也拿不准到底该贴哪个标签。这也再次印证了，现阶段的内容审核不能完全指望人工智能，最后一步肯定还得让人工来拍板。


\chapter{基于预训练模型的用户评论多标签智能审核系统的设计与实现}
\section{系统需求与总体架构设计}
为了将前文训练好的算法模型真正在业务中跑通，本章介绍如何开发一套完整的可视化审核系统。在动手写代码之前，我们先梳理了系统的核心需求：在功能上，系统必须支持审核员手敲单条文本测试，也得支持上传表格进行批量处理；同时，因为模型是个“黑盒”，系统需要把各项违规的概率和命中的敏感词直观地展示出来。在性能上，最基本的要求就是推理响应要快，界面在频繁点击时不能卡死。

结合这些需求，我们在软件架构上采用了四层解耦的设计模式，这样不仅代码写起来更清晰，以后如果要换更强的大模型也方便替换。各层的具体分工如下：

\begin{enumerate}
	\def\labelenumi{\arabic{enumi}.}
	\item
	\textbf{数据交互层}：这是系统对外的窗口。主要负责接收用户的两种输入——要么是网页文本框里敲的一句话，要么是包含成千上万条历史评论的 CSV 文件。
	\item
	\textbf{模型推断层}：相当于系统的“大脑”。它负责把训练好的 RoBERTa-wwm-ext 权重加载到显存里。当收到上层传来的文本后，它会快速跑完分词和矩阵运算，吐出 6 个违规标签的原始概率。
	\item
	\textbf{业务逻辑与规则层}：这里写了我们自己定制的过滤规则。它会用正则去抓手机号，用词典去抓脏话，算出惩罚分后，再跟模型吐出的概率进行融合打分。最后根据系统设定的“及格线”来判定这句话到底违不违规。
	\item
	\textbf{前端应用层}：我们使用 Streamlit 框架写了一个 Web 界面，把底层那些枯燥的数字变成了漂亮的柱状图和警报弹窗，并且会把每次审核的结果悄悄记录到后台日志里。
\end{enumerate}

\section{核心功能模块实现}
\subsection{单条与批量预测模块}
“单条预测”主要是给审核员做实时抽查用的。当用户在输入框打完字点击提交后，系统会在零点几秒内跑完模型推理和规则匹配。接着，界面上会画出一个柱状图，把这句话在色情、辱骂等 6 个维度上的危险程度清清楚楚地标出来。如果哪一项超标了，就会直接打上高危标签，提醒审核员注意。

而“批量预测”则是为了应付那种堆积如山的历史数据。用户只要把待查的 CSV 文件拖进网页，系统后台就会利用 PyTorch 的 \texttt{DataLoader} 机制，把数据按批次打包塞进 GPU 狂跑。跑完之后，系统会自动吐出一份新的 CSV 表格。表里不仅保留了原来的文字，还多出了各项概率、命中的规则以及最终的“死刑判决书”，极大地解放了审核专员的双手。

\subsection{阈值动态调节模块}
在实际的审核环境里，审核的标准从来不是一成不变的。有时候碰上突发的热点事件，某类违规内容就会激增。为了应付这种情况，我们在网页侧边栏放了一排滑动条。

最上面的是“全局阈值”，用来控制整体审核的松紧。下面则是针对 6 个标签的独立阈值。假如最近网上因为某件事爆发了严重的地域黑，管理员只要把“地域歧视”那个滑块从 0.5 往左拉到 0.3。这样一来，系统对地域黑的容忍度就会瞬间降低，能抓出更多阴阳怪气的疑似违规评论交给人去复核，非常灵活。

\subsection{规则提示与日志审计模块}
由于光给出一个“违规”的结论很难服众，我们在界面上加了一个“命中规则展示”的模块。只要一句话踩了雷，界面马上会用高亮红字提示“触发了某个敏感词”或者“检测到发微信号引流”，这就让整个黑盒模型变得有理有据了。

此外，为了方便出了问题能追责，系统会默默把每一次点击审核的记录写进后端的日志文件里。日志里存了当时的时间、那句评论说了啥、各项概率是多少，以及最终是怎么判的。这批日志后续不仅能用来抓 Bug，还能抽出来当做下一代模型重新训练的语料。

\section{核心模块详细设计与性能优化}
\subsection{模型单例加载与显存优化机制}
在开发过程中我们遇到了一个很头疼的问题：预训练模型的参数量上亿，每次初始化加载权重都要好几秒。如果用户在网页上随便拖一下滑动条或者切个页面，Streamlit 默认会把整个代码重新跑一遍，导致模型被反复加载，不仅卡得要死，还很容易把显存撑爆。

为了解决这个大坑，我们引入了 Streamlit 的 \texttt{@st.cache\_resource} 装饰器。这个语法的神奇之处在于，它能帮我们实现\textbf{模型和分词器的全局单例缓存}。也就是说，只有在系统刚启动的第一次，模型会被老老实实加载进显存里；后面不管是哪个用户、点什么按钮发来的预测请求，系统都会直接复用内存里已经躺好的那个模型实例。这样改完之后，单条文本的检测延迟直接从几秒钟降到了 50 到 100 毫秒左右，完全感觉不到卡顿。

\subsection{批量审核 DataLoader 批处理机制}
在处理批量 CSV 数据时，如果我们用普通的 \texttt{for} 循环一条一条去喂给模型，不仅速度慢如蜗牛，而且根本发挥不出显卡的并行计算优势。

所以我们在批量预测模块里重写了底层的逻辑，自己继承了 PyTorch 的 \texttt{Dataset} 类，并把它塞进了 \texttt{DataLoader} 里。程序会把上传的 DataFrame 切成一个个大小为 32（\texttt{batch\_size=32}）的数据包，然后再统一搬到 GPU 显存里去做矩阵乘法。算完之后再把结果拼回原表里。等所有进度条跑完，前端页面会弹出一个 \texttt{st.download\_button} 的下载按钮，提醒用户带打分结果的报告已经生成好了。

\section{系统界面与交互流程设计}
整个系统页面被干脆利落地切成了两块：\textbf{侧边控制台}和\textbf{主工作区}。

\begin{enumerate}
	\def\labelenumi{\arabic{enumi}.}
	\item
	\textbf{侧边控制台}：
		页面左侧就是前面提到的动态滑动条面板。只要审核员更改改了某个阈值，内存里的判断标准立刻随之变化，下一个进来的请求马上按新阈值进行判断。
\begin{figure}[!htb]
	\centering
	\includegraphics[width=0.5\textwidth]{picture51}
	\caption{侧边控制台}
	\label{figureLabel}
\end{figure}

	
		
	\item
	\textbf{主工作区交互（以单条审查为例）}：
	主工作区上面是一个菜单栏，支持在单条预测和批量预测这两大个大功能之间来回切。
		用户输入文本一敲回车，页面右上角马上开始转圈圈，提示模型正在疯狂计算。
		结果正常的话，如果语句无明显违规，则提示说“该内容未发现明显违规”；如果查出违规内容，则提示“该内容存在违规风险，建议审核或屏蔽”。无论是正常还是违规的评论，系统都会利用图表组件画出柱状图并给出判断结果，直白地显示这句话命中了哪个违规条例。如果命中规则库里的关键词，则同时显示语句里的哪个关键词违规了，让算法的判决令人心服口服。
	\begin{figure}[!htb]
	\centering
	\includegraphics[width=1\textwidth]{PIC52}
	\caption{违规评论测试}
	\label{figureLabel}
\end{figure}
\begin{figure}[!htb]
	\centering
	\includegraphics[width=1\textwidth]{PIC53}
	\caption{正常评论测试}
	\label{figureLabel}
\end{figure}
\end{enumerate}


\chapter{总结与展望}
\section{全文总结}
本课题主要是为了解当今各大论坛和社交平台里评论区的网络空间安全隐患问题。以前那种纯靠人工或者只盯着敏感词去匹配的审核方式，要么处理速度太慢，要么总是有漏网之鱼。为了改变这个现状，我们引入了深度学习里的预训练大模型技术，基于 RoBERTa-wwm-ext 开发了一个能同时给一条评论打上多种违规标签的智能审核系统。

回顾整个项目的推进过程，我们主要完成了以下几项工作。首先是搞定了最基础的数据问题，我们清洗并整理出了一份包含五万多条中文评论的数据集，并重点解决了里面“色情辱骂多、冷门歧视词汇少”的数据不平衡问题。接着，我们在代码层面搭建了多标签分类网络，选用了更合适的 Sigmoid 激活与 BCE 损失函数。最终，经过调参微调，模型在测试集上跑出了超过 0.99 的 Macro-F1 分数，分类效果比较理想。

当然，光有一个深度学习模型还不够。为了防止这个“黑盒”模型偶尔犯傻，漏掉一些极为明显的违规网址或脏话，我们还在系统里硬编码了一套传统规则引擎来做兜底，实现了“模型概率+规则惩罚”的双重过滤。最后，为了让这套算法真能落地使用，我们用 Streamlit 框架写了个可视化的交互页面，把单条预测、批量刷数据、拖拽调阈值以及记录审核日志这些功能全都缝合了进去。总的来说，这套系统基本达到了“机器先秒级初筛，拿不准的再交给人看”的设计初衷，对于中小型的社区网站来说，具备一定的直接应用价值。

\section{未来展望}
虽说目前这套系统已经能顺畅运行，各项测试指标也达到了初期的设想，但互联网上的网络黑话和变体词更新速度实在是太快了，对于网络空间安全仍然构成挑战。如果以后有机会继续深入研究，我觉得系统还有以下几个地方可以好好改进：

\begin{enumerate}
	\def\labelenumi{\arabic{enumi}.}
	\item
	\textbf{完善黑话词典与正常语境防误伤}：模型现在偶尔还会误伤正常讨论（比如把医学科普或者正规新闻里的词当成涉黄违规）。以后可以添加一个动态更新的白名单，结合更深度的上下文消歧技术，尽量减少这种“老实人被误封”的现象。
	\item
	\textbf{借助生成式大模型搞定冷门数据}：像一些很小众的方言骂人或者刚流行起来的谐音梗，由于我们的训练集里根本没见过，目前模型是认不出来的。未来可以考虑接一个像 Deepseek 这样的大语言模型，让它们帮忙生成一些对抗样本来扩充我们的训练集，从而治好模型的“偏科”问题。
	\item
	\textbf{建立人工纠错的反馈闭环}：现在的模型训好之后，参数权重就焊死了，这其实不太符合实际的审核业务。理想的状态应该是：审核员在后台把机器判错的评论手工改过来后，系统能把这些“错题”悄悄攒在一起，每隔一段时间自动在后台触发一次微调训练，让模型越用越聪明。
	\item
	\textbf{从纯文字走向多模态审核}：目前这套代码只能看纯文字的评论。但现在网友发评论都喜欢带表情包或者各种动图。以后要是能接入 OCR（光学字符识别）和图像分类技术，把涉黄图片和违规表情包也一并管理起来，那这套审核系统才算真正完整。
\end{enumerate}




\chapterspecialcenter{致 谢}
夏天开始的故事，也终于要在夏天迎来结尾。完成这篇论文后，我也将在22岁的年纪完成我儿时一直憧憬的本科生涯。
落其实者，思其树。我要感谢我的毕设指导老师刘玲君老师，在毕业设计完成的过程中给予了我很大的发挥空间的同时，也用严谨的态度对我的论文进行指导和修正。
父母半生恩，兄弟一世情。我要感谢家人，家人们用最大的力量将我托举到不属于我的高度，也给了我很多力量是我一直以来最坚实的后盾。从幼儿园到大学，离家的路越来越远，不变是你们一直都在我身边，你们的温暖让我一直坚信只有家人才是神赐福的完美艺术，唯有不断的向前才能不辜负你们的期待。
好酒易求，挚友难遇。我要感谢生命中的挚友们，你们是我生命组成的一部分，见证了我的幸福和失落，以真诚一直延续着缘分。挚友是亲自挑选的家人，愿我们都能在自己选择的道路上一帆风顺。
道阻且长，行则将至。最后我要感谢我自己，你脸上带着微笑，心中却充满悲苦，无论如何你都坚持到了现在，何时葡萄才成熟，你要静候再静候。你不是英雄，而是幸存者，安慰自己这一切是生命的被迫扩容。终有一天你会走出这段雨季，希望你能看清自己，接受自己，活成自己。



\end{document}
